%%%%%%%%%%%%%%%%%%%%%%%%%%%%%%%%%%%%%%%%%%%%%%%%%%%%%%%%%%%%%%%%%%%%%%%%%%%%%%%
%% Internal Reproduction Report — Dawes+2023 (ApJS 269:61) autocorrelation
%% search for multiply lensed and binary quasars in the DESI Legacy Surveys.
%%
%% Single-column `article` class with the shared tech-report preamble at
%% reproductions/tech-report.sty. Builds with vanilla TeX Live (>= 2022).
%%%%%%%%%%%%%%%%%%%%%%%%%%%%%%%%%%%%%%%%%%%%%%%%%%%%%%%%%%%%%%%%%%%%%%%%%%%%%%%
\documentclass[11pt,letterpaper]{article}
\usepackage{../../tech-report}

% Dawes-specific math macros (\asec is already provided by amsmath; we
% define a robust as-text arcsecond marker under a private name instead)
\newcommand{\asec}{\ensuremath{^{\prime\prime}}}
\newcommand{\PMSIG}{\ensuremath{\mathrm{PMSIG}}}
\newcommand{\PXSIG}{\ensuremath{\mathrm{PXSIG}}}
\newcommand{\degsq}{\ensuremath{\mathrm{deg}^2}}

\title{Reproducing the Dawes 2022/2023 Autocorrelation Search for\\
       Multiply Lensed and Binary Quasars in the DESI Legacy Surveys}
\author{Greg Benson \and Xiaosheng Huang}
\date{June 2026}
\addbibresource{references.bib}

\begin{document}

\maketitle
\subtitle{Internal Reproduction Report}

\begin{abstract}
We reproduce the algorithmic core of the \citet{dawes2023quasars}
``autocorrelation'' search for multiply lensed and binary quasars in the
DESI Legacy Imaging Surveys. The published autocorrelation step is a
friends-of-friends (FoF) angular grouping of quasar targets at a fixed
linking length --- searched first at $10\asec$, then halved to ${<}5\asec$
--- followed by Gaia EDR3 proper-motion and parallax cuts and human visual
grading. Because Dawes' ${\sim}5{\rm M}$-target \emph{photometric} DESI Quasar
Sample is not on disk, we run the identical \texttt{spherimatch} FoF
\autocite{spherimatch} on the most directly available proxy --- the
$1{,}645{,}843$-spectrum DESI DR1 \emph{spectroscopic} QSO sample from
\texttt{zall-pix-iron.fits}. The proxy is ${\sim}1/3$ the size and is
selected differently, so absolute candidate counts cannot match the
published $436$; what we reproduce faithfully is the \emph{algorithm} and
the \emph{recovery of discoverable systems}. Our headline result: every
published candidate that is \emph{discoverable in our proxy} (has $\geq 2$
distinct DR1 spectroscopic QSOs within $5\asec$) is correctly grouped by
our FoF --- $58/58 = 100\,\%$, the apples-to-apples analog of the paper's
$94/94 = 100\,\%$ recovery of its ``discoverable known'' systems. Recovered
matches sit a median $0.139\asec$ from the published centroid. The full
$1.6{\rm M}$-QSO FoF runs in ${\sim}3$\,s. Visual-inspection grading
(Grades A/B/C) is explicitly out of scope, exactly as in the companion
\citet{hsu2025pairwise} reproduction.
\end{abstract}

\tableofcontents

\bigskip

% =============================================================================
\section{Introduction}
% =============================================================================

This report documents the reproduction of \citet{dawes2023quasars}
(ApJS~269:61; arXiv submitted 2022, hence the ``Dawes 2022/2023'' label),
one of a series of internal reproductions covering
Prof. Xiaosheng Huang's strong-gravitational-lensing program. Companion
reports cover \citet{huang2020decals} (DECaLS ResNet lens finder) and
\citet{hsu2025pairwise} (pair-wise spectroscopic lens search over DESI
DR1); the present work is the quasar-lens analog of the latter.

\paragraph{What the paper did.} Multiply lensed quasars are powerful
cosmological probes: the time delay between their multiple images yields
an independent measurement of $H_0$
\autocite{treu2016marshall,dawes2023quasars}. To grow the small known
sample, \citet{dawes2023quasars} search the $19{,}000\,\degsq$ DESI Legacy
Imaging Surveys \autocite{dey2019legacysurveys} for new lensed and binary
quasars. They apply an \emph{autocorrelation algorithm} to the
${\sim}5$~million objects classified as quasars in the DESI Quasar Sample,
visually inspect and rank the resulting systems, and present
\textbf{436 new multiply lensed and binary quasar candidates}.

\paragraph{The autocorrelation algorithm (\S3).} The paper's
``autocorrelation'' step groups quasar targets that lie within an angular
separation cut and ``recursively connects multiple pairs of quasar targets
within the same candidate system (such as would occur in `triples' and
`quads')''. This is precisely a \textbf{friends-of-friends} grouping with a
fixed angular linking length. Dawes first search a $10\asec$ radius
(``to err on the side of completeness''), then halve the cut to ${<}5\asec$
because ${>}95\,\%$ of the discoverable known systems have images separated
by ${<}5\asec$. They test the algorithm on the 94 ``discoverable known''
systems and on simulations, \textbf{achieving a 100\,\% success rate}, then
deploy it on the full ${\sim}5{\rm M}$ quasar-target sample.

\paragraph{Lensed-quasar system bookkeeping (Table~1).} Of the
${\sim}300$ known lensed-quasar systems in the footprint, $94$ are
``discoverable known'' (have $\geq 2$ objects in the DESI Quasar Sample so
the algorithm \emph{can} group them), of which $72$ are ``discoverable
confirmed''. The new output is the $436$ candidates.

\paragraph{Post-FoF stellar rejection (\S4.1).} Milky-Way stellar
contaminants are rejected using Gaia EDR3 \autocite{gaiaedr3} proper motion
and parallax. Dawes define significances
$\PMSIG = \sqrt{\mu_\alpha^2 + \mu_\delta^2}/\sqrt{\sigma_{\mu_\alpha}^2 +
\sigma_{\mu_\delta}^2}$ and $\PXSIG = \varpi/\sigma_\varpi$, and accept a
source as extragalactic if $\PMSIG < 8$ \emph{and} $\PXSIG < 3.5$. Gaia
information is available for at least one image of $380/436 \approx 87\,\%$
of the published candidates.

\paragraph{Scope.} As with the \citet{hsu2025pairwise} reproduction, the
published pipeline has an algorithmic stage (autocorrelation FoF + Gaia
cuts) followed by a human visual-inspection / grading stage (Grades
A/B/C). \textbf{This reproduction covers the algorithmic stage only.}
Visual grading is not attempted.

% =============================================================================
\section{The proxy sample (and why absolute counts cannot match)}
% =============================================================================

The single most important caveat of this reproduction is the input
catalog. Dawes ran on the DESI Legacy Imaging Surveys \emph{photometric}
DESI Quasar Sample --- ${\sim}5$~million Random-Forest-selected QSO
\emph{targets}, one detection per source. \textbf{That photometric target
catalog is not on disk.} We therefore substitute the most directly
available proxy: the DESI DR1 \emph{spectroscopic} QSO sample carved from
the shared $28.4{\rm M}$-fiber \texttt{zall-pix-iron.fits}
\autocite{desidr1} via the pre-filter cascade in Table~\ref{tab:prefilter}.

\begin{table}[ht]
\centering
\caption{QSO pre-filter cascade on \texttt{zall-pix-iron.fits} (DR1).}
\label{tab:prefilter}
\begin{tabular}{lr}
\toprule
filter                                     & remaining \\
\midrule
raw rows                                   & 28{,}425{,}963 \\
\texttt{ZCAT\_PRIMARY == True}             & 27{,}547{,}223 \\
\texttt{ZWARN == 0}                        & 20{,}415{,}199 \\
\texttt{SPECTYPE == QSO}                   & \textbf{1{,}645{,}843} \\
\midrule
Dawes input (photometric, NOT used)        & ${\sim}5{,}000{,}000$ \\
\bottomrule
\end{tabular}
\end{table}

This proxy is a related but \textbf{not identical} sample: it is
${\sim}1/3$ the size and selected differently (spectroscopically confirmed
QSOs only, vs.\ photometric RF-selected targets). Two consequences follow
and are quantified honestly throughout this report:
\begin{enumerate}
\item \textbf{Absolute candidate counts cannot match $436$.} Most published
      candidate \emph{images} are faint, blended lensed images that DESI
      never spectroscopically confirmed as QSOs. Only $138/436$ published
      candidates have \emph{any} DR1 QSO within $5\asec$, and only
      $58/436$ have $\geq 2$ \emph{distinct} DR1 QSOs (i.e.\ are
      ``discoverable'' in our proxy).
\item \textbf{The reproduction targets are the algorithm and the recovery
      of discoverable systems}, not the raw $436$. We assess these in
      \S\ref{sec:recovery}.
\end{enumerate}

% =============================================================================
\section{Pipeline}
% =============================================================================

The pipeline is implemented in seven numbered scripts at
\texttt{reproductions/dawes-2022/}, summarized in Table~\ref{tab:scripts}.
The FoF engine is the exact \texttt{spherimatch} \autocite{spherimatch}
tool used by the \citet{hsu2025pairwise} pipeline; the catalog fetch and
Gaia queries use \texttt{astroquery} \autocite{astroquery2019}; sky-coordinate
matching uses \texttt{astropy} \autocite{astropy2022}.

\begin{table}[ht]
\centering
\caption{Reproduction scripts.}
\label{tab:scripts}
\small
\begin{tabular}{lp{0.66\linewidth}}
\toprule
script & purpose \\
\midrule
\texttt{00\_fetch\_published\_catalog.py} &
  download Dawes Table~2 from VizieR \texttt{J/ApJS/269/61}
  ($875$ image-rows $\to 436$ systems). \\
\texttt{01\_build\_qso\_sample\_and\_fof.py} &
  build the DR1 QSO proxy sample; \texttt{spherimatch} FoF at $5\asec$
  and $10\asec$. \\
\texttt{02\_recover\_published\_candidates.py} &
  nearest-group recovery of the $436$ published positions. \\
\texttt{03\_gaia\_pm\_px\_cuts.py} &
  Gaia EDR3 $\PMSIG$/$\PXSIG$ stellar rejection (checkpointed/resumable). \\
\texttt{04\_make\_figures.py} &
  separation histogram + recovery-offset histogram. \\
\texttt{05\_refine\_candidates.py} &
  collapse spectroscopic duplicates, keep resolved
  ($0.3\asec \le \mathrm{sep} \le 5\asec$) candidate systems. \\
\texttt{06\_recovery\_ceiling.py} &
  proxy-aware \emph{conditional} recovery (analog of $94/94$). \\
\bottomrule
\end{tabular}
\end{table}

\paragraph{FoF grouping.} For each of the two linking lengths
($\ell = 5\asec$ and $10\asec$) we call \texttt{spherimatch.fof} on the
(\texttt{TARGET\_RA}, \texttt{TARGET\_DEC}) of the $1{,}645{,}843$ QSOs and
retain groups of size $\geq 2$. The $5\asec$ run takes $2.8$\,s, the
$10\asec$ run $2.8$\,s; loading the $28.4{\rm M}$-row catalog dominates the
wall-clock at $30$\,s.

\paragraph{Refinement.} The raw FoF groups are dominated by pairs at
${\sim}0\asec$ separation: distinct \texttt{TARGETID}s at the \emph{same}
sky position (the same physical QSO observed under multiple DESI
surveys/programs, or Tractor splits). Dawes' photometric target catalog has
one detection per source and so never sees this pile-up. Script~05 collapses
near-zero-separation duplicates (dedup within $0.2\asec$) and keeps only
``resolved'' systems with maximum internal image separation in
$[0.3\asec, 5\asec]$ --- $0.3\asec$ matches the minimum published Dawes
separation of $0.32\asec$.

% =============================================================================
\section{Algorithmic-step results}
% =============================================================================

The headline comparison of our output against \citet{dawes2023quasars} is
in Table~\ref{tab:headline}. Because the input samples differ in size and
selection, the absolute group counts are expected to differ and are shown
for transparency rather than as a match claim; the rows that constitute the
actual reproduction target (recovery of discoverable systems) are in
\S\ref{sec:recovery}.

\begin{table}[ht]
\centering
\caption{Algorithmic-step output: published vs ours. The
$\dagger$ rows compare counts across \emph{different input samples}
(${\sim}5{\rm M}$ photometric targets vs $1.6{\rm M}$ spectroscopic QSOs)
and are not expected to match.}
\label{tab:headline}
\small
\begin{tabular}{lrr}
\toprule
metric & published (Dawes) & ours (DR1 QSO proxy) \\
\midrule
QSO input sample                          & ${\sim}5{,}000{,}000$ phot. & 1{,}645{,}843 spec. \\
FoF groups @ $5\asec$ (raw)$^\dagger$    & ---                   & 10{,}917 \\
FoF groups @ $10\asec$ (raw)$^\dagger$   & ---                   & 12{,}896 \\
recommendations @ $10\asec$$^\dagger$    & ${>}27{,}000$ targets & 25{,}962 spectra \\
recommendations @ $5\asec$$^\dagger$     & ${\sim}6{,}000$ targets & 21{,}965 spectra \\
resolved candidates @ $5\asec$$^\dagger$ & 436 candidates        & 1{,}548 \\
\bottomrule
\end{tabular}
\end{table}

\paragraph{Multiplicity breakdown.} The raw $5\asec$ FoF groups split into
$10{,}790$ doubles, $123$ triples and $4$ quads ($\geq 4$ images); the
$10\asec$ run gives $12{,}730/162/4$. After refinement (\S5) the
resolved systems are $1{,}529$ doubles, $18$ triples and $1$ quad, a
double-to-multi ratio of ${\approx}80{:}1$. Dawes report
${\approx}21{:}1$ in the text ($24$ quads identified by the algorithm; the
published Type column lists $432$ Double $+\,4$ Quad). The ratio cannot
match because spectroscopic confirmation of \emph{all} faint images of a
quad is rare in DR1 --- most published quads were Tractor-deblended from
photometric data we have no access to.

\paragraph{Separation distribution.} Figure~\ref{fig:sep} contrasts the
maximum within-group image separation of our $5\asec$ FoF groups against
the published Dawes separations. Our distribution piles up at $0\asec$
(the spectroscopic-duplicate artifact); the published separations span
$0.32$--$4.99\asec$ with median $1.9\asec$, the physically meaningful
regime that script~05's refinement isolates.

\begin{figure}[ht]
\centering
\includegraphics[width=\linewidth]{figures/separation_hist.png}
\caption{Maximum within-group image separation: our $5\asec$ FoF groups
(filled, $N=10{,}917$) vs the published Dawes+2023 separations
(red outline, $N=436$). The $0\asec$ spike in our distribution is distinct
\texttt{TARGETID}s at the same sky position (the same physical QSO observed
under multiple DESI surveys/programs); the refinement step (\S5) removes it.}
\label{fig:sep}
\end{figure}

% =============================================================================
\section{Recovery of the published candidates}
\label{sec:recovery}
% =============================================================================

The paper's headline target is \textbf{100\,\% recovery of the 94
``discoverable known'' systems}. We assess recovery at two levels.

\paragraph{Raw recovery (all 436).} Cross-matching each of the $436$
published candidate positions against our FoF group centroids
(recovered if the nearest centroid is within $3\asec$) gives
\textbf{$61/436 = 14.0\,\%$}, identical at $5\asec$ and $10\asec$, with
a median offset of $\mathbf{0.139\asec}$ for recovered systems. The
by-grade breakdown is $17/102$ (A), $15/118$ (B), $29/216$ (C). This raw
fraction is \emph{capped by the proxy}: only $138/436$ candidates have any
DR1 QSO within $5\asec$ at all, so $\sim\!86\,\%$ of the misses are simply
candidates with no spectroscopic counterpart to group. Figure~\ref{fig:off}
shows the bimodal offset distribution --- a tight recovered population at
${\sim}0.1\asec$ and a large unrecoverable population whose nearest QSO
group is ${\sim}30\asec$ away.

\begin{figure}[ht]
\centering
\includegraphics[width=\linewidth]{figures/recovery_offsets.png}
\caption{Offset from each of the $436$ published candidates to the nearest
$5\asec$ FoF group centroid. $61$ fall within $3\asec$ (dashed line);
the recovered population sits at median $0.139\asec$. The
${\sim}30\asec$ pile-up is published candidates with no nearby DR1
spectroscopic QSO --- a proxy limitation, not an algorithmic failure.}
\label{fig:off}
\end{figure}

\paragraph{Conditional recovery --- the apples-to-apples analog of 94/94.}
The fair comparison restricts to candidates that are \emph{discoverable in
our proxy}: those with $\geq 2$ distinct DR1 spectroscopic QSOs within
$5\asec$ (the proxy analog of Dawes' ``$\geq 2$ objects in the DESI Quasar
Sample'' definition of ``discoverable known''). There are $58$ such
candidates, and our FoF algorithm correctly groups
\textbf{all $58$ of them: $58/58 = 100.0\,\%$}. This is the direct analog of
the paper's $94/94 = 100\,\%$ recovery, and it reproduces.

\begin{table}[ht]
\centering
\caption{Recovery results. The conditional row is the apples-to-apples
analog of the published $94/94$.}
\label{tab:recovery}
\begin{tabular}{lrr}
\toprule
quantity & published (Dawes) & ours (DR1 QSO proxy) \\
\midrule
discoverable systems                  & 94             & 58 \\
\;of which grouped by FoF             & 94 ($100\,\%$) & \textbf{58 ($\mathbf{100\,\%}$)} \\
raw recovery of all candidates        & n/a            & $61/436 = 14.0\,\%$ \\
median offset of recovered            & ---            & $0.139\asec$ \\
candidates with any DR1 QSO ${<}5\asec$ & n/a        & $138/436$ \\
\bottomrule
\end{tabular}
\end{table}

% =============================================================================
\section{Gaia EDR3 stellar rejection}
% =============================================================================

Script~03 reproduces the Dawes \S4.1 Gaia EDR3 \autocite{gaiaedr3}
contaminant cut. For every QSO image in the $1{,}548$ resolved candidate
groups ($3{,}179$ images) it performs a $2\asec$ cone search of
\texttt{gaiaedr3.gaia\_source} via \texttt{astroquery.gaia}
\autocite{astroquery2019}, computes $\PMSIG$ and $\PXSIG$ from the published
formulae, and flags a group as stellar if any matched image has
$\PMSIG \geq 8$ or $\PXSIG \geq 3.5$. The query is per-position and
checkpointed (the bulk TAP upload-join hangs on the host), so it is the
slowest step.

\paragraph{Status.} Because faint lensed images mostly fall below Gaia's
${\sim}21$\,mag limit, a large fraction of our QSO images have \emph{no}
Gaia match --- that absence is itself part of the quasar signature, and
Dawes report Gaia information for $\geq 1$ image of only $380/436
\approx 87\,\%$ of candidates. At the time of this report the per-position
Gaia crawl over all $3{,}179$ images was still in progress (the ESA Gaia
archive is throttled in the run-up to DR4 and times out intermittently);
the cut logic is implemented and validated on the queried subset, but the
final per-group survival count is not yet finalized and is therefore
reported as an incomplete step rather than a headline number.

% =============================================================================
\section{Reproducibility caveats / not reproduced}
% =============================================================================

\paragraph{Proxy sample (the dominant caveat).} Every absolute-count
discrepancy in this report traces to the input substitution: a
$1.6{\rm M}$-spectrum DR1 spectroscopic QSO sample standing in for Dawes'
${\sim}5{\rm M}$ photometric RF-selected target catalog, which is not on
disk. The proxy is ${\sim}1/3$ the size and selected differently. We make no
claim of matching the $436$ candidate count, the
${>}27{,}000$/${\sim}6{,}000$ recommendation counts, or the $\approx 21{:}1$
double-to-quad ratio; these compare quantities derived from different input
samples.

\paragraph{Spectroscopic-duplicate pile-up.} The raw FoF groups are
dominated by ${\sim}0\asec$ pairs (distinct \texttt{TARGETID}s at the same
sky position). This is an artifact of using a multi-survey
\emph{spectroscopic} catalog and is absent from Dawes' one-detection-per-source
photometric catalog. Script~05's dedup-and-resolve step removes it, yielding
the $1{,}548$ resolved systems.

\paragraph{Double:quad ratio.} Our $80{:}1$ (1{,}529 doubles, 1 quad) vs
the published ${\approx}21{:}1$ ($24$ quads). Spectroscopic confirmation of
\emph{all} images of a quad is rare in DR1; most published quads were
Tractor-deblended from photometric imaging we cannot access.

\paragraph{Visual grading.} The reduction from algorithmic candidates to
the graded A/B/C catalog is a human visual-inspection workflow, out of
scope here exactly as in the \citet{hsu2025pairwise} reproduction.

\paragraph{Gaia crawl incomplete.} As noted in \S7, the full per-image Gaia
EDR3 query had not finished at report time; the survival count is pending.

% =============================================================================
\section{Software and data availability}
% =============================================================================

\begin{itemize}
\item Scripts: \texttt{reproductions/dawes-2022/00--06\_*.py}; this report
      builds against the artifacts they produce in \texttt{data/} and
      \texttt{figs/}.
\item Published catalog: VizieR \texttt{J/ApJS/269/61} table~2
      (\citealt{dawes2023quasars}; $875$ image-rows, $436$ systems), pulled
      via \texttt{astroquery.vizier} \autocite{astroquery2019}.
\item DESI DR1 zcatalog source (shared with the \citet{hsu2025pairwise}
      reproduction):
      \url{https://data.desi.lbl.gov/public/dr1/spectro/redux/iron/zcatalog/v1/zall-pix-iron.fits}
      (22.4~GB, public).
\item \texttt{spherimatch} \autocite{spherimatch}: PyPI, MIT licensed; the
      same FoF tool used by the \citet{hsu2025pairwise} pipeline.
\item Gaia EDR3 \autocite{gaiaedr3} via \texttt{astroquery.gaia}
      \autocite{astroquery2019}; sky matching via \texttt{astropy}
      \autocite{astropy2022}.
\item Python environment: \texttt{/home/benson/.venvs/hsu} ---
      \texttt{spherimatch}, \texttt{astropy}, \texttt{astroquery},
      \texttt{pyarrow}, \texttt{pandas}, \texttt{pdfminer.six}.
\item Code repository: \url{https://github.com/usfcs-edu/agentic-lensing}.
\end{itemize}

% =============================================================================
\section{Conclusions}
% =============================================================================

The autocorrelation algorithm at the heart of \citet{dawes2023quasars} is a
friends-of-friends angular grouping; reproduced with the exact
\texttt{spherimatch} FoF tool at the published $5\asec$ and $10\asec$
linking lengths, it \textbf{recovers $100\,\%$ of the systems discoverable
in our sample ($58/58$)} --- the apples-to-apples analog of the paper's
$94/94 = 100\,\%$ recovery of its discoverable-known systems --- with a
median centroid offset of $0.139\asec$. The raw $14\,\%$ recovery and the
mismatches in absolute candidate counts and double-to-quad ratio are driven
\emph{entirely} by the proxy input sample (DR1 spectroscopic QSOs,
${\sim}1/3$ the size, differently selected), not by the algorithm. The full
$1.6{\rm M}$-QSO FoF runs in ${\sim}3$~seconds on local hardware. The Gaia
EDR3 stellar-rejection step is implemented per the paper's $\PMSIG < 8$,
$\PXSIG < 3.5$ definitions but its per-position crawl was incomplete at
report time, and visual grading remains out of scope.

\bibliographystyle{plainnat}
\bibliography{references}

\end{document}
